\documentclass[11pt,a4paper]{article}

\usepackage[utf8]{inputenc}
\usepackage[T1]{fontenc}
\usepackage[margin=1in]{geometry}
\usepackage{amsmath, amssymb, amsthm, mathtools}
\usepackage{hyperref}
\usepackage{xcolor}
\usepackage[expansion=false]{microtype}
\usepackage{enumitem}
\usepackage{booktabs}

\hypersetup{
  colorlinks=true,
  linkcolor=black,
  urlcolor=blue!50!black,
  citecolor=black,
}

\newtheorem{theorem}{Theorem}
\newtheorem{corollary}[theorem]{Corollary}
\newtheorem{lemma}[theorem]{Lemma}
\newtheorem{observation}[theorem]{Observation}
\newtheorem*{residue*}{Residue}
\newtheorem*{remark*}{Remark}

\setlength{\parskip}{0.4em}

\title{A Strange Idea, Followed Further\\
\large Step 6.3, Actually Done, and What That Showed}
\author{Vladyslav Savytskyy\\ \small with Claude (Anthropic)}
\date{May 2026, Argentina \\ \emph{companion to ``A Strange Idea, Followed Honestly''}}

\begin{document}

\maketitle

\begin{abstract}
The previous document stopped one section short of doing the calculation it
described. This one does the calculation. The result is in every respect
consistent with the standard result of Chamseddine--Connes--van Suijlekom (2013).
Along the way, one parenthetical statement of the previous document was found to
be wrong as written and is corrected here; one numerical claim was redone with
current data; and the open question that the previous document called Step~4 is
located one click more precisely. The strange idea did not become any less
strange. It became better mapped.
\end{abstract}

{\setlength{\parskip}{0pt}\tableofcontents}

\bigskip

\section{What this document is, and why it exists}

The previous document\footnote{V. Savytskyy with Claude (Anthropic), \emph{A Strange
Idea, Followed Honestly: The Generator Hypothesis After Every Stress Test We
Could Apply}, May 2026.} closed with Section~6.3, in which we had attempted to
set up the explicit $(\widehat A, J, \gamma)$-bimodule for the Pati--Salam
algebra in code, verified the KO-dimension~6 conditions, and then stopped. The
stated reasons for stopping were:
\begin{enumerate}[label=(\arabic*),leftmargin=2em]
  \item Reproducing the published calculation adds no new information.
  \item The genuine open question is not answered by reproducing the calculation.
\end{enumerate}
Both reasons are defensible. Neither, on reflection, is sufficient. The
calculation that was paraphrased in three lines is also the calculation that, if
one wants to test \emph{anything} downstream of the order-one carve-out, one
needs to actually have running. The previous document had a parenthetical
assertion about $A_F$ failing to embed as a unital $*$-subalgebra of the
Pati--Salam algebra; that assertion is, under a literal reading, false, and it
was made because we had not yet built the embedding. The previous document also
described the Koide residual to seven decimal places using a number whose
provenance it did not carefully chase; that number is also approximately right
but is scheme-dependent, and the provenance turns out to matter slightly.

This document records the calculation done. It is, like its predecessor, not a
theory and not a framework. It is a record. The audience is presumed to be
familiar with the previous document; the section numbering continues from there.

\section{Recap, in two paragraphs}

The Generator hypothesis was a maximalist claim that the Standard Model and
gravitation are consequences of a single self-referential algebraic object. Most
of its alleged empirical content turned out to be fitted rather than derived; a
structural rigidity theorem ruled out its central object under its own axioms;
the honest rewrite reduced the framework to operator-algebraic renormalization
group theory, which is real, well-studied, and predicts nothing new. The
residual mathematical question, on which the framework's chance of becoming
non-trivial entirely rests, is the selection of the finite algebra
$A_F = \mathbb{C} \oplus \mathbb{H} \oplus M_3(\mathbb{C})$ in Connes's spectral
action program. This is Step~4. As of May 2026 nobody has closed it.

The previous document located Step~4 inside the algebra-enumeration step of
Chamseddine--Connes (2007) and at the level of the bimodule decoration. It did
not, however, build the bimodule. The present document does.

\section{The Koide residual, redone carefully}\label{sec:koide}

Before the structural work, one minor numerical cleanup. The previous document
quoted
\[
  K = \frac{m_e + m_\mu + m_\tau}{(\sqrt{m_e}+\sqrt{m_\mu}+\sqrt{m_\tau})^2}
    = 0.666661 \pm 0.000007,
\]
deviation from $2/3$ at the $10^{-5}$ level. The provenance of the value was
not specified, and the uncertainty was given without specifying the input set.

Recomputed in this session against current data:
\begin{itemize}[leftmargin=2em]
  \item With CODATA~2022 values for $m_e$, $m_\mu$, $m_\tau$ (the last from
        PDG~2022 propagated into CODATA): $K = 0.6666605(1)$, deviation from
        $2/3$ of $\sim 9 \times 10^{-6}$, statistical distance from $2/3$
        approximately $0.91\sigma$ under the quoted uncertainties.
  \item With PDG~2024 values directly, and in particular the updated
        $m_\tau = 1776.93 \pm 0.09$~MeV: $K = 0.6666645(5)$, deviation from
        $2/3$ of $\sim 2 \times 10^{-6}$, statistical distance approximately
        $0.43\sigma$.
\end{itemize}

\begin{observation}\label{obs:koide-current}
For exact agreement $K = 2/3$, the tau mass would need to shift by between
$+0.04$ and $+0.11$~MeV, well inside the current PDG uncertainty.
\end{observation}

\begin{observation}\label{obs:koide-scheme}
$K$ is computed from pole masses. Under one-loop pure QED running it is
RG-invariant; under full Standard Model running it is not. The
``$10^{-5}$~level coincidence'' is therefore a coincidence in a specific
renormalization scheme. Under a different scheme the deviation can be order
$10^{-3}$ or larger.
\end{observation}

This does not refute the relation. It does mean that the previous document's
report of seven-decimal agreement was, on close inspection, agreement in a
scheme that the document did not name. The honest version is: $K$ is consistent
with $2/3$ within current experimental uncertainty, in the pole-mass scheme;
the deviation is at most $1\sigma$, possibly less; no derivation of why this
should hold exists in any rigorous framework. The relation does not extend
cleanly to quarks. It may be coincidence; it may be a hint; it is currently
undecidable. This is the same verdict the previous document reached, with one
fewer false-precision figure.

\section{The embedding, actually constructed}\label{sec:embedding}

The previous document, in Observation~9, stated:

\begin{quote}\itshape
Of the four CC07-irreducible candidates at fermion dimension~32 ($M_4(\mathbb{C})$,
$M_4(\mathbb{R}) \oplus M_4(\mathbb{R})$, $M_2(\mathbb{H}) \oplus M_4(\mathbb{R})$,
$M_2(\mathbb{H}) \oplus M_2(\mathbb{H})$), none contains $A_F$ as a unital
$*$-subalgebra. Even the canonical Pati--Salam algebra $M_2(\mathbb{H}) \oplus
M_4(\mathbb{C})$ does not.
\end{quote}

The parenthetical about $M_2(\mathbb{H}) \oplus M_4(\mathbb{C})$ is, taken at
face value, wrong. An explicit unital $*$-embedding exists and we constructed it.

\begin{theorem}\label{thm:embedding}
There exists a unital $*$-algebra homomorphism
\[
  \iota \colon \mathbb{C} \oplus \mathbb{H} \oplus M_3(\mathbb{C})
  \;\longrightarrow\; M_2(\mathbb{H}) \oplus M_4(\mathbb{C})
\]
with zero kernel.
\end{theorem}

\begin{proof}
Explicit construction. View $\mathbb{C} \subset \mathbb{H}$ as the subalgebra
spanned by $\{1, i\}$. For $(\lambda, q, X) \in \mathbb{C} \oplus \mathbb{H}
\oplus M_3(\mathbb{C})$, define
\[
  \iota(\lambda, q, X) \;=\;
  \Big(\, \operatorname{diag}(\lambda, q) \,\in\, M_2(\mathbb{H}),
        \;\; \operatorname{diag}(\lambda, X) \,\in\, M_4(\mathbb{C}) \,\Big),
\]
where on the right $\operatorname{diag}(\lambda, X)$ places $\lambda$ in the
$(1,1)$ entry and $X$ in the lower $3 \times 3$ block. Unitality, additivity,
$\mathbb{R}$-linearity, $*$-preservation and faithfulness all follow by direct
check. Multiplicativity reduces to multiplicativity of the inclusion
$\mathbb{C} \hookrightarrow \mathbb{H}$ and of block-diagonal matrix
multiplication. The image is unital because both blocks are. The kernel is zero
because each summand of the source is faithfully represented in at least one
summand of the target.
\end{proof}

We verified this numerically as a sanity check on the embedding. Over fifty
independent random elements of $A_F$:
\begin{itemize}[leftmargin=2em]
  \item unitality residual: $0$;
  \item multiplicativity residual: $\le 1.78 \times 10^{-15}$
        (machine epsilon, double precision);
  \item additivity, $\mathbb{R}$-linearity, $*$-preservation: residuals all $0$;
  \item faithfulness check (nonzero elements mapping to zero):
        $0$ out of $100$.
\end{itemize}

This is not a deep correction. The previous document's Observation~9 was
attempting to express a real and well-known fact in the NCG literature, namely
that $A_F$ is not a \emph{canonical multiplicity-1} unital $*$-subalgebra of the
Pati--Salam algebra. The K-rank obstruction it pointed to ($10 \neq 12$) is
real. But the parenthetical sentence, as written, slid from ``no canonical
inclusion'' to ``no inclusion'', which is a different and stronger and false
claim. The right statement is the one Chamseddine--Connes--van Suijlekom~2013
make: $A_F$ is carved out of the Pati--Salam algebra at the level of the
$(J,\gamma)$-decorated bimodule, by the order-one condition acting on a fixed
grading, not by the existence or non-existence of an inclusion of algebras at
the matrix level. Both relationships hold; one is much more informative than
the other; the previous document was making a sharp point and overshooting
slightly. The replacement statement is:

\begin{observation}\label{obs:embedding-corrected}
$A_F$ embeds as a unital $*$-subalgebra of $M_2(\mathbb{H}) \oplus M_4(\mathbb{C})$
via a multiplicity-2 / corner-3 inclusion. It does not embed as a canonical
multiplicity-1 unital $*$-subalgebra: the K-rank sum of $A_F$'s natural module
($2 + 4 + 4 = 10$) does not equal that of the Pati--Salam natural module
($8 + 4 = 12$). The order-one carve-out from the Pati--Salam algebra to $A_F$
is therefore a statement about the bimodule, not the algebra.
\end{observation}

This is what Section~6.3 of the previous document was getting at, before it
stopped. The next four sections do not stop.

\section{The bimodule, actually built}\label{sec:bimodule}

For one fermion generation, set
\[
  H_F \;=\; \mathbb{C}^{32}
       \;=\; \operatorname{span} \big\{\,|a, s, w, c\rangle \;:\;
              a \in \{\pm 1\},\; s \in \{\pm 1\},\;
              w \in \{\pm 1\},\; c \in \{1,2,3,4\}\,\big\},
\]
with $a$ indexing particle ($+$) vs antiparticle ($-$), $s$ left ($+$) vs right
($-$), $w$ the up/down member of the weak doublet, and $c$ the colour, with
$c \in \{1,2,3\}$ quark colours and $c = 4$ the lepton ``colour.''

\paragraph{Grading and real structure.}
Define
\[
  \gamma\, |a, s, w, c\rangle \;=\; (a \cdot s)\, |a, s, w, c\rangle,
  \qquad
  J\, |a, s, w, c\rangle \;=\; |\!-\!a, s, w, c\rangle
  \;\text{ followed by complex conjugation of coefficients.}
\]
With these choices: $\gamma^2 = +1$, $J^2 = +1$, $J\gamma + \gamma J = 0$.
KO-dimension~6 holds on the nose. The dimension count matches:
$\operatorname{tr}\gamma = 0$, i.e.\ equal numbers of $+$ and $-$ chirality states.

\paragraph{Algebra action.}
For $(\lambda, q, m) \in A_F$, define $\pi_F(\lambda, q, m) \in \operatorname{End}(H_F)$
as follows. On the particle sector ($a = +1$):
\begin{itemize}[leftmargin=2em]
  \item Left-handed sector ($s = +1$): the quaternion $q$ acts on the
        $(w = +1, w = -1)$ doublet via its $M_2(\mathbb{C})$ realisation
        $q \mapsto \begin{pmatrix} q_0 + i q_1 & q_2 + i q_3 \\
        -q_2 + i q_3 & q_0 - i q_1 \end{pmatrix}$, with $c$ trivial.
  \item Right-handed up sector ($s = -1$, $w = +1$): scalar multiplication
        by $\lambda$.
  \item Right-handed down sector ($s = -1$, $w = -1$): scalar multiplication
        by $\bar\lambda$.
\end{itemize}
On the antiparticle sector ($a = -1$): $m \in M_3(\mathbb{C})$ acts on the
quark colour index $c \in \{1,2,3\}$, and $c = 4$ is fixed (identity action).
The opposite action is $\pi_F^{\mathrm{op}}(a) = J\, \pi_F(a)^*\, J^{-1}$.

\paragraph{Verification.}
All of the following were checked at machine precision in this session. We
record the worst-case residuals over the random samples or basis elements used,
in double-precision arithmetic.

\begin{center}
\renewcommand{\arraystretch}{1.1}
\begin{tabular}{lll}
\toprule
condition & meaning & residual \\
\midrule
$\pi_F(1) = I$ & $\pi_F$ unital & $0$ \\
$\pi_F(ab) = \pi_F(a)\pi_F(b)$ & $\pi_F$ multiplicative & $1.83 \times 10^{-15}$ \\
$\pi_F(a^*) = \pi_F(a)^*$ & $*$-preserving & $0$ \\
$[\gamma, \pi_F(a)] = 0$ & algebra preserves grading & $0$ \\
$[\pi_F(a), \pi_F^{\mathrm{op}}(b)] = 0$ & bimodule property & $0$ \\
$J^2 = +I$ & KO-dim 6 sign of $J$ & $0$ \\
$J\gamma + \gamma J = 0$ & KO-dim 6 anti-commutation & $0$ \\
$\gamma^2 = I$ & $\mathbb{Z}/2$ grading & $0$ \\
$D = D^*$ & Dirac is Hermitian & $0$ \\
$D\gamma + \gamma D = 0$ & Dirac is odd & $0$ \\
$DJ = JD$ & Dirac commutes with $J$ & $0$ \\
$[[D, \pi_F(a)], \pi_F^{\mathrm{op}}(b)] = 0$ & order-one for $A_F$ & $0$ \\
\bottomrule
\end{tabular}
\end{center}

\bigskip

The non-zero residual $1.83 \times 10^{-15}$ for multiplicativity is the
unavoidable round-off in composing $32 \times 32$ complex matrix products in
double precision; one cannot expect better. The other entries are exact zeros
because the corresponding conditions are computed entry-by-entry from
permutations and additions, not from matrix multiplications. In short: the
linear algebra was done correctly.

\paragraph{The Dirac operator.}
Define $D \in \operatorname{End}(H_F)$ on basis vectors by, for each
$(w, \mathrm{type})$ where $\mathrm{type} \in \{\mathrm{quark}, \mathrm{lepton}\}$,
choosing a Yukawa value $y_{w, \mathrm{type}}$ and setting
\[
  D\, |+1, +1, w, c\rangle \;=\; y_{w, \mathrm{type}}\, |+1, -1, w, c\rangle,
  \qquad
  D\, |+1, -1, w, c\rangle \;=\; \overline{y_{w, \mathrm{type}}}\, |+1, +1, w, c\rangle,
\]
with the antiparticle entries fixed by $DJ = JD$. The construction makes $D$
Hermitian, anti-commuting with $\gamma$, and commuting with $J$.

\paragraph{Order-one.}
For this $D$ and for $A_F$:
\[
  \big[\![ D, \pi_F(a) ], \pi_F^{\mathrm{op}}(b) \big]
  \;=\; 0
  \qquad \text{for all } a, b \in A_F,
\]
verified at machine precision over random samples.

\paragraph{The carve-out, on the nose.}
The natural representation of $M_2(\mathbb{H}) \oplus M_4(\mathbb{C})$ acting
on the same $H_F$ extends $\pi_F$ in the canonical way: $M_2(\mathbb{H})$ acts
on the $(s,w)$ block of particles (treating $(L_u, L_d, R_u, R_d)$ as
$\mathbb{H}^2$), and $M_4(\mathbb{C})$ acts on the colour index of
antiparticles including $c = 4$. Two things are then computationally manifest.
\begin{itemize}[leftmargin=2em]
  \item The off-diagonal quaternionic blocks of $M_2(\mathbb{H})$ mix $s = +1$
        with $s = -1$, so $[\gamma, \pi_{PS}(X)] \neq 0$ for generic
        $X \in M_2(\mathbb{H}) \oplus M_4(\mathbb{C})$. Residual under
        random sampling: $\sim 8.2$, which is to say not zero.
  \item Restricting to the $\gamma$-preserving subalgebra
        $\mathbb{H}_L \oplus \mathbb{H}_R \oplus M_4(\mathbb{C})$, the
        order-one condition for the SM Dirac fails:
        residual $\sim 9.6 \times 10^{-2}$.
\end{itemize}

\begin{observation}\label{obs:carveout}
With the bimodule above, $A_F$ satisfies the order-one condition for the SM
Yukawa Dirac. The natural $\gamma$-preserving extension
$\mathbb{H}_L \oplus \mathbb{H}_R \oplus M_4(\mathbb{C})$ does not. The
non-$\gamma$-preserving extension $M_2(\mathbb{H}) \oplus M_4(\mathbb{C})$ does
not even satisfy $[\gamma, \pi] = 0$. The carve-out is, in this concrete
computational realisation, a fact about the bimodule.
\end{observation}

This is the calculation that the previous document said one would have to do.
We have now done it. The conclusion is the standard one. Section~6.3 of the
previous document is, with this section, closed in the modest sense of having
the calculation it pointed at on the table.

\section{The Dirac-dimension functional}\label{sec:diracdim}

Once the bimodule is built, one can ask a question that the previous document
did not ask, because it had not yet built the bimodule: for a subalgebra
$\mathcal{A} \subseteq \operatorname{End}(H_F)$ acting on $H_F$ compatibly with
$\gamma$ and $J$, how large is the space of admissible Dirac operators?

Concretely, define
\[
  \mathcal{D}(\mathcal{A}) \;:=\; \Big\{\, D \in \operatorname{End}(H_F)
  \;:\;
  D = D^*,\;
  D\gamma + \gamma D = 0,\;
  DJ = JD,\;
  [[D, \pi(a)], \pi^{\mathrm{op}}(b)] = 0\;\forall a, b \in \mathcal{A}
  \,\Big\},
\]
a real linear subspace of $\operatorname{End}(H_F)$.

The first three conditions (Hermiticity, anti-commutation with $\gamma$,
commutation with $J$) carve a base subspace of $\operatorname{End}(H_F)$. With
$\dim_{\mathbb{R}} \operatorname{End}(H_F) = 2 \cdot 32^2 = 2048$ and the
conditions linearly independent, direct computation gives the base null space
dimension $d_{\mathrm{base}} = 272$. The order-one condition then cuts this
$272$-dimensional space further, by an amount depending on $\mathcal{A}$.

We computed $\dim_{\mathbb{R}} \mathcal{D}(\mathcal{A})$ for a chain of
subalgebras. Setting up the constraints in a basis of the base null space and
taking ranks via QR with column pivoting, the result is the following.

\begin{center}
\renewcommand{\arraystretch}{1.05}
\begin{tabular}{lrr}
\toprule
Algebra $\mathcal{A}$ & $\dim_{\mathbb{R}} \mathcal{A}$ &
$\dim_{\mathbb{R}} \mathcal{D}(\mathcal{A})$ \\
\midrule
none (i.e.\ only the unit) & 1 & 272 \\
$\mathbb{C}$ & 3 & 146 \\
$\mathbb{H}$ & 5 & 146 \\
$\mathbb{C} \oplus \mathbb{H}$ & 7 & 80 \\
$M_3(\mathbb{C})$ & 19 & 128 \\
$\mathbb{C} \oplus M_3(\mathbb{C})$ & 21 & \textbf{16} \\
$\mathbb{H} \oplus M_3(\mathbb{C})$ & 23 & \textbf{16} \\
$A_F = \mathbb{C} \oplus \mathbb{H} \oplus M_3(\mathbb{C})$ & 25 & \textbf{16} \\
$\mathbb{H}_L \oplus \mathbb{H}_R \oplus M_4(\mathbb{C})$ & 41 & 8 \\
\bottomrule
\end{tabular}
\end{center}

\bigskip

Two things are worth noting before any interpretation.

\begin{observation}\label{obs:nonmonotone}
$\dim \mathcal{D}$ is \emph{not} a monotone function of $\dim \mathcal{A}$.
$M_3(\mathbb{C})$ ($\dim 19$) gives $\dim \mathcal{D} = 128$, while the strictly
smaller $\mathbb{C} \oplus \mathbb{H}$ ($\dim 7$) gives $\dim \mathcal{D} = 80$.
The functional depends on which directions of $\operatorname{End}(H_F)$ the
algebra hits, not on the algebra's dimension.
\end{observation}

\begin{observation}\label{obs:plateau}
Three different subalgebras of $A_F$ give $\dim \mathcal{D} = 16$: the algebras
$\mathbb{C} \oplus M_3(\mathbb{C})$, $\mathbb{H} \oplus M_3(\mathbb{C})$, and
$A_F$ itself. The full algebra $A_F$ does not impose any further constraint on
$\mathcal{D}$ beyond what either $\mathbb{C} \oplus M_3(\mathbb{C})$ or
$\mathbb{H} \oplus M_3(\mathbb{C})$ already imposes.
\end{observation}

The number $16$ is not arbitrary. Per generation, the SM Yukawa parameters are
encoded as a $2 \times 2$ Yukawa matrix per fermion type (quark, lepton),
each entry a complex number. Modulo Hermiticity and $J$-pairing, this gives
$2 \times 4 \times 2 = 16$ real parameters per generation. The plateau at $16$
is exactly the per-generation SM Yukawa parameter count.

This is checkable. Set up $16$ explicit physical Yukawa basis Diracs, one per
$(\mathrm{type}, w_L, w_R, \mathrm{re}/\mathrm{im})$. Each one was verified to
lie in $\mathcal{D}(A_F)$ with residual $\le 1.6 \times 10^{-14}$. Their image
in the $16$-dimensional coordinate space of $\mathcal{D}(A_F)$ has rank $16$.
So:

\begin{theorem}\label{thm:exactly-yukawa}
$\mathcal{D}(A_F)$ has dimension $16$, and is spanned exactly by the
$16$-parameter family of physical Yukawa Diracs (per generation, modulo
Hermiticity and $J$-pairing).
\end{theorem}

This is what Chamseddine--Connes--van Suijlekom~2013 say. With the bimodule on
the table, we have just checked it.

\section{$A_F$ is not picked out by $\dim \mathcal{D}$ alone}\label{sec:nonunique}

Observation~\ref{obs:plateau} says the Dirac-dimension functional, taken by
itself, does not single out $A_F$ among its subalgebras. We pushed on this and
verified the strongest version of the statement.

\begin{theorem}\label{thm:subspaces-equal}
As subspaces of $\operatorname{End}(H_F)$,
\[
  \mathcal{D}\big(\mathbb{C} \oplus M_3(\mathbb{C})\big)
  \;=\;
  \mathcal{D}\big(\mathbb{H} \oplus M_3(\mathbb{C})\big)
  \;=\;
  \mathcal{D}(A_F).
\]
\end{theorem}

\begin{proof}
The first two subalgebras are contained in $A_F$, so each imposes a subset of
$A_F$'s constraints, hence each $\mathcal{D}$ \emph{contains} $\mathcal{D}(A_F)$.
By Theorem~\ref{thm:exactly-yukawa} and Observation~\ref{obs:plateau} all three
spaces are 16-dimensional. Containment of finite-dimensional spaces of equal
dimension is equality.

Numerical confirmation: we computed orthonormal bases for all three subspaces
and projected one onto another. The projection residuals were
$\le 7.2 \times 10^{-15}$ in all three pairwise comparisons.
\end{proof}

This sharpens, in a way that is initially surprising and on reflection slightly
embarrassing, what the Dirac-dimension functional can and cannot do.

\begin{observation}\label{obs:redundancy}
The $\mathbb{H}$ summand of $A_F$ imposes zero additional order-one constraint
on $\mathcal{D}$ beyond what is already imposed by $\mathbb{C} \oplus M_3(\mathbb{C})$.
Symmetrically, the $\mathbb{C}$ summand imposes zero additional constraint
beyond what $\mathbb{H} \oplus M_3(\mathbb{C})$ already imposes. The two
summands $\mathbb{C}$ and $\mathbb{H}$ each independently suffice to cut
$\mathcal{D}(M_3(\mathbb{C}))$ ($\dim 128$) down to the SM Yukawa subspace
($\dim 16$). Adding the other contributes nothing further at the level of
$\mathcal{D}$.
\end{observation}

This is a real structural fact. It also means that any honest variational
attempt to pick $A_F$ out using $\dim \mathcal{D}$ as the functional will fail.
The maximum is attained on a plateau, and the plateau contains $A_F$ as one
point among at least three. Any unique selection of $A_F$ must use additional
information.

The natural additional information is the gauge group $U(\mathcal{A})$. We
record the relevant dimensions and note the salient gap.
\begin{itemize}[leftmargin=2em]
  \item $\dim U(\mathbb{C} \oplus M_3(\mathbb{C})) = 1 + 9 = 10$, with
        $U(\mathbb{C} \oplus M_3(\mathbb{C})) = U(1) \times U(3)$. No
        $SU(2)$ factor.
  \item $\dim U(\mathbb{H} \oplus M_3(\mathbb{C})) = 3 + 9 = 12$, with
        $U(\mathbb{H} \oplus M_3(\mathbb{C})) = SU(2) \times U(3)$.
  \item $\dim U(A_F) = 1 + 3 + 9 = 13$, with
        $U(A_F) = U(1) \times SU(2) \times U(3)$, reducing modulo
        $\mathbb{Z}_6$ to the SM gauge group.
\end{itemize}
The two non-$A_F$ algebras at $\dim \mathcal{D} = 16$ are inadequate for the
following different reasons.
\begin{itemize}[leftmargin=2em]
  \item $\mathbb{C} \oplus M_3(\mathbb{C})$: no $SU(2)_L$. No $W$ bosons.
  \item $\mathbb{H} \oplus M_3(\mathbb{C})$: has $SU(2)$, has $U(1)$, but the
        $U(1)$ is the determinant phase of $U(3)$ and the $\lambda \leftrightarrow
        \bar\lambda$ distinction on right-handed up vs.\ down fermions is
        absent. The Standard Model hypercharge assignments
        $Y(u_R) \neq Y(d_R)$ require the action of the $\mathbb{C}$ summand on
        the bimodule, where one gets $\lambda$ on $u_R$ and $\bar\lambda$ on
        $d_R$. Without that summand the hypercharge assignments are wrong.
\end{itemize}

Putting these together:

\begin{observation}\label{obs:two-conditions}
$A_F$ is the minimal $*$-subalgebra of the natural $\gamma$-preserving extension
of the SM bimodule such that
\begin{itemize}[leftmargin=2em,nosep]
  \item $\dim_{\mathbb{R}} \mathcal{D}(\mathcal{A}) \ge 16$ (full SM Yukawa
        freedom), \emph{and}
  \item the unitary group $U(\mathcal{A})$ contains both the $SU(2)_L$ factor
        and the $U(1)$ factor that distinguishes $Y(u_R)$ from $Y(d_R)$ in the
        bimodule action.
\end{itemize}
\end{observation}

Two conditions. Not one.

\section{The 8-dimensional Pati--Salam Dirac space}\label{sec:ps8dim}

The carve-out from Pati--Salam to $A_F$, then, is from $\dim \mathcal{D}
= 8$ to $\dim \mathcal{D} = 16$. The factor of two suggests a clean structural
identification, and indeed there is one.

For the algebra $\mathbb{H}_L \oplus \mathbb{H}_R \oplus M_4(\mathbb{C})$
(which preserves $\gamma$) acting on the same bimodule via the canonical
extension of $\pi_F$, we computed $\mathcal{D}$ as an 8-dimensional subspace
of $\mathcal{D}(A_F)$, and expressed each of its basis vectors in the
physical-Yukawa coordinate system of Theorem~\ref{thm:exactly-yukawa}.

The output is, in every case, of the following form. A basis vector of the
8-dimensional surviving Dirac space, when written in physical Yukawa
coordinates, has the structure
\[
  \sum_{j} c_j \cdot (\mathrm{quark\ Yukawa})_j
  \;+\;
  \sum_{j} c_j \cdot (\mathrm{lepton\ Yukawa})_j,
\]
\emph{with the same coefficient $c_j$ on the quark term and on the corresponding
lepton term.} Each of the 16 individual Yukawa basis directions (8 quark,
8 lepton) violates the additional constraints by an $O(1)$ amount. Only
quark--lepton paired combinations survive.

\begin{theorem}\label{thm:ps-leptoquark}
The eight surviving Dirac directions in $\mathcal{D}(\mathbb{H}_L \oplus
\mathbb{H}_R \oplus M_4(\mathbb{C}))$ are exactly the leptoquark-unified
combinations: a combined $(\mathrm{quark} + \mathrm{lepton})$ Yukawa for each
of the 8 $(w_L, w_R, \mathrm{re}/\mathrm{im})$ assignments, with equal weight
on quark and lepton.
\end{theorem}

This is, in plain language, the Pati--Salam prediction that leptons are the
fourth colour, hence that lepton and quark Yukawas are tied at the unification
scale. Numerically, after RG running, this is the source of the
$m_b/m_\tau \approx 3$ low-scale prediction. The carve-out from
$\mathbb{H}_L \oplus \mathbb{H}_R \oplus M_4(\mathbb{C})$ to $A_F$ is the
breaking of this tie: the order-one condition for the $\mathbb{C}$ summand of
$A_F$ acting only on right-handed Standard Model particles, together with the
restriction from $M_4(\mathbb{C})$ to $M_3(\mathbb{C}) \oplus (\text{trivial on
leptons})$, is precisely what allows independent quark and lepton Yukawas.

We did not, in this session, observe anything in this analysis that is not
already in Chamseddine--Connes--van Suijlekom~2013. What is new is the
computational receipt: the 8-dimensional space is, as a subspace of the
16-dimensional one, exactly the diagonal quark--lepton matched subspace.
We have not seen this stated in this form in the published literature.
We have also not searched the published literature exhaustively. So either
this is a known consequence stated in slightly different language elsewhere,
or it has been computed before without being highlighted. Both possibilities
are likely; both possibilities are fine; the receipts are in
\texttt{phase6\_8dim.py}.

\section{Step 4, restated}\label{sec:step4-sharpened}

The previous document stated the open question this way:

\begin{quote}\itshape
The variational principle that would close the framework, if it exists, is a
functional on $\mathcal{M} = \{\,\text{Krajewski diagrams of KO-dim~6}\,\}
\times \{\,\text{compatible Dirac moduli}\,\}$ whose minimum (or canonical
critical point) is the SM bimodule.
\end{quote}

This is still the question. We can now say slightly more about what such a
functional must do.

\begin{observation}\label{obs:step4-shape}
A functional $F$ on $\mathcal{M}$ that singles out $A_F$ uniquely cannot
factor through $\dim \mathcal{D}$ alone, because that functional plateaus.
It must combine, in some genuine way, at least the following two kinds of
information:
\begin{itemize}[leftmargin=2em,nosep]
  \item information about the bimodule structure (KO-dimension, order-one
        carve-out, dimension of admissible Dirac operators);
  \item information about the unitary group of $\mathcal{A}$ acting on the
        bimodule (the gauge content, including hypercharge assignments).
\end{itemize}
A merely topological or merely representation-theoretic functional will not
suffice; the function must see both kinds of structure simultaneously. The
candidate spectral-action functional has this property in principle, since the
heat-kernel coefficients $a_2, a_4$ involve both eigenvalue content of $D$
and representation content of $\mathcal{A}$. Whether the unique minimizer is
$A_F$ remains, as of May 2026, open.
\end{observation}

The previous document's framing of Step~4 stands. We add one observation: the
naive simplification ``maximize the dimension of the admissible Dirac space''
does not work, and the reason it does not work is now visible at the level of
explicit subspaces. A successful single functional must do more than count.

\section{Residues, after this round}\label{sec:residues-2}

The residues from the previous document remain residues. The only one we
touched is Koide, and its status is unchanged in substance.

\begin{residue*}[Koide, current status.]
$K = 0.6666645(5)$ in the pole-mass scheme using PDG~2024 data; consistent
with $2/3$ to within $0.43\sigma$; not preserved under full SM running. A
tau-mass shift of $+0.04$ to $+0.11$~MeV (inside current uncertainty)
would make it exact. No derivation; no quark extension. Coincidence or hint,
currently undecidable. Per Observation~\ref{obs:koide-scheme}, the
seven-decimal version of the statement is scheme-specific.
\end{residue*}

\begin{residue*}[Three generations.]
No change. Empirical fact, no principled derivation in any rigorous
framework.
\end{residue*}

\begin{residue*}[Cosmological constant.]
No change. Worst quantitative discrepancy in theoretical physics; addressed by
no current framework.
\end{residue*}

\begin{residue*}[Selection of $A_F$.]
Sharper after Sections~\ref{sec:diracdim}--\ref{sec:step4-sharpened}.
Not closed.
The cleanest available characterisation of $A_F$ now goes: it is the minimal
$*$-subalgebra such that (i) the order-one condition on the SM bimodule admits
the full 16-dimensional Yukawa space, and (ii) the unitary group action carries
the SM gauge group with correct hypercharge assignments. This is a two-condition
characterisation, not a one-functional minimisation. Closing Step~4 still
requires the latter.
\end{residue*}

\begin{residue*}[Unreasonable effectiveness of operator algebras.]
No change. The alignment between purely internal operator-algebraic structures
and physical observers, measurements, scales, and effective theories remains
unexplained. We did not propose an explanation.
\end{residue*}

\begin{residue*}[The meta-puzzle, after Wigner.]
No change. Still 1960.
\end{residue*}

\section{Verdict, updated}\label{sec:verdict-2}

The previous document's verdict stands without subtraction. We add three lines.

The Generator hypothesis, as originally proposed, did not survive honest
scrutiny. That has not changed.

The honest rewrite (categorical RG on operator algebras) is mathematical
physics in a particular dialect and is not new. That has not changed.

The closest existing program (Connes's spectral action) recovers the bosonic
Standard Model from one functional given $A_F$. That has not changed.

What changed in this session:
\begin{itemize}[leftmargin=2em]
  \item The Section~6.3 calculation that the previous document gestured at is
        now on the table, end-to-end, at machine precision.
  \item One parenthetical sentence of the previous document
        (Observation~\ref{obs:embedding-corrected}) was corrected.
  \item The functional approach to Step~4 was tested at the level of
        $\dim \mathcal{D}$, and shown to plateau in a way that rules out the
        naive single-functional minimisation. The successful functional must
        see both bimodule structure and gauge content, simultaneously.
  \item The structural meaning of the 8-vs-16 dimensional carve-out was
        located: it is precisely the breaking of Pati--Salam leptoquark
        unification at the level of Yukawa parameters.
\end{itemize}

None of these closes Step~4. The previous document's summary remains:

\begin{quote}\itshape
We did not close it. We did locate it.
\end{quote}

We located it slightly more precisely. That is the substantive content of this
document, and it is smaller than the calculation that produced it would
suggest. This is normal. Honest work has this shape.

\section{What this document is}\label{sec:what-this-is}

This document is, like its predecessor, a record. It is not a theory, not a
framework, not a derivation of the universe from one equation. The arithmetic
in it was done in Python in a sandboxed environment with a wall clock; the
linear algebra ran in BLAS; the SVDs were computed with LAPACK; the
constructions matched the published literature wherever they overlapped.
Receipts for everything in Sections~\ref{sec:embedding}--\ref{sec:ps8dim} are
in the files \texttt{phase2a\_explicit\_embedding.py},
\texttt{phase2c\_carveout.py}, \texttt{phase3\_dirac\_dim\_v2.py},
\texttt{phase4\_fast.py}, \texttt{phase5\_structure.py}, and
\texttt{phase6\_8dim.py}. The numerical precision of every reported residual
can be reproduced by anyone with a Python installation; the patience required
is approximately one hour of compute on a 4~GB single-core machine.

The strange idea, in the form it survives after this further round, is
exactly as strange as it was before, and not strange in any new way. It still
points at one open question. It has not become a theory; it has not become a
framework. The thing it did, in this round, was to acquire one more layer of
precision around the open question. The next round, if there is one, would
either close the question or determine that it cannot be closed without
genuinely new mathematical ideas.

Both possibilities remain interesting. Neither is currently known. The
strangeness has not gone anywhere. It has been measured a little more
carefully.

That is enough.

\appendix

\section{Computational receipts}\label{app:receipts}

This section records, for each substantive computational claim in the body of
the document, the file in which it was carried out, the principal verification
performed, and the worst-case numerical residual obtained. All computations
were performed in Python with \texttt{numpy} 2.4.4 and \texttt{scipy} 1.17.1 in
double precision, in a sandboxed environment with no external network access.
The wall-clock cost of reproducing the full set is approximately one hour on a
single-core 4~GB machine.

\begin{center}
\renewcommand{\arraystretch}{1.15}
\begin{tabular}{p{4.6cm}p{6cm}p{3.4cm}}
\toprule
file & verifies & worst residual \\
\midrule
\texttt{phase1\char`_koide.py} &
Koide ratio under CODATA 2022 and PDG 2024; tau-mass shift for $K = 2/3$;
$\sigma$-distance &
data-driven (see Sec.~\ref{sec:koide}) \\
\texttt{phase2a\char`_v2.py} &
enumeration of 9421 finite real semisimple $*$-algebras up to real dim 96
and 5 summands; CC07-admissible count 91; unital $*$-subalgebra obstruction
for $A_F$ in fermion-dim-32 candidates &
combinatorial \\
\texttt{phase2a\char`_explicit\char`_\allowbreak embedding.py} &
the explicit map $\iota: A_F \to M_2(\mathbb{H}) \oplus M_4(\mathbb{C})$ of
Theorem~\ref{thm:embedding}; unitality, multiplicativity, additivity,
$\mathbb{R}$-linearity, $*$-preservation, faithfulness over 50 random elements
each &
$1.78 \times 10^{-15}$ \\
\texttt{phase2c\char`_carveout.py} &
the full bimodule of Section~\ref{sec:bimodule}; all entries of the
verification table; order-one for $A_F$; the $[\gamma, \pi_{PS}] \neq 0$
obstruction; order-one failure for
$\mathbb{H}_L \oplus \mathbb{H}_R \oplus M_4(\mathbb{C})$ &
$1.83 \times 10^{-15}$ for non-zero entries; obstruction residuals $8.23$
and $9.63 \times 10^{-2}$ \\
\texttt{phase3\char`_dirac\char`_\allowbreak dim\char`_v2.py} &
$d_{\mathrm{base}} = 272$ via SVD of the stacked
Hermitian+anti-$\gamma$+commute-$J$ constraint matrix;
$\dim \mathcal{D}(A_F) = 16$;
$\dim \mathcal{D}(\mathbb{H}_L \oplus \mathbb{H}_R \oplus M_4(\mathbb{C})) = 8$
via QR-based rank determination &
SVD tolerance $10^{-10} \cdot \sigma_{\max}$ \\
\texttt{phase4\char`_fast.py} &
the full chain of subalgebras in Section~\ref{sec:diracdim} using batched
\texttt{numpy} order-one constraint computation; the plateau
$\dim \mathcal{D} = 16$ for $\mathbb{C} \oplus M_3$,
$\mathbb{H} \oplus M_3$, and $A_F$; identification of the 16 physical Yukawa
basis Diracs and rank check &
$\le 1.64 \times 10^{-14}$ \\
\texttt{phase5\char`_structure.py} &
subspace equality
$\mathcal{D}(\mathbb{C} \oplus M_3) = \mathcal{D}(\mathbb{H} \oplus M_3) =
\mathcal{D}(A_F)$ via projector residual; containment in
$\mathcal{D}(M_3)$; projection of the explicit SM Yukawa Dirac onto
$\mathcal{D}(A_F)$; rank verification of the physical Yukawa basis &
$\le 7.20 \times 10^{-15}$ for subspace equality;
$6.69 \times 10^{-17}$ for SM Dirac projection \\
\texttt{phase6\char`_8dim.py} &
the 8-dimensional Dirac space of
$\mathbb{H}_L \oplus \mathbb{H}_R \oplus M_4(\mathbb{C})$ computed by projecting
its additional order-one constraints onto $\mathcal{D}(A_F)$; expression of
each surviving basis direction in the physical Yukawa coordinate system;
verification that quark and lepton coefficients in each surviving direction
are paired (equal) &
matched-coefficient pairing exact within tolerance $10^{-8}$ \\
\bottomrule
\end{tabular}
\end{center}

\bigskip

The python files are not included here; they exist in the working environment
in which this document was produced. Anyone wishing to reproduce the
computations should be able to write equivalent scripts in less than a day
from the specifications in Sections~\ref{sec:embedding} through
\ref{sec:ps8dim} of this document.

\section{Notes on what a successful Step~4 functional must do}\label{app:step4-shape}

We do not, in this document or elsewhere, propose a functional that closes
Step~4. Observation~\ref{obs:step4-shape} restricts the shape of any such
functional. We record here a slightly longer list of properties such a
functional $F$ would need to have, in case it is useful for anyone attempting
to construct one.

\begin{enumerate}[leftmargin=2em]
  \item $F$ must depend on the bimodule $(A, H_F, \gamma, J, \pi)$, not just on
        the algebra $A$. The Dirac-dimension functional already has this
        property; the plateau at $\dim \mathcal{D} = 16$ across three different
        subalgebras of $A_F$ shows that bimodule structure alone is not
        sufficient either.
  \item $F$ must distinguish $A_F$ from $\mathbb{H} \oplus M_3(\mathbb{C})$
        despite the two having the same Dirac space.
        The visible discriminator is the action of the $\mathbb{C}$ summand of
        $A_F$ on the right-handed up vs.\ down sectors via $\lambda$ vs.\
        $\bar\lambda$, which is what gives the $Y(u_R) \neq Y(d_R)$ structure.
        Any successful $F$ must therefore see hypercharge assignments at the
        level of the representation, not just at the level of the gauge group.
  \item $F$ must distinguish $A_F$ from $\mathbb{C} \oplus M_3(\mathbb{C})$
        despite the same Dirac space. Here the visible discriminator is the
        presence of $SU(2)_L$ in the gauge group. Any successful $F$ must
        therefore involve a non-abelian piece.
  \item $F$ must, taken across the chain of Section~\ref{sec:diracdim}, single
        out $A_F$ over $\mathbb{H}_L \oplus \mathbb{H}_R \oplus M_4(\mathbb{C})$,
        i.e.\ choose the SM over Pati--Salam. This requires
        $F$ to prefer the un-unified Yukawa structure
        ($\dim \mathcal{D} = 16$) over the leptoquark-unified structure
        ($\dim \mathcal{D} = 8$). On the face of it, this looks like preferring
        \emph{more} Dirac freedom, which is a sensible criterion. Any
        successful $F$ must, however, also choose $A_F$ specifically over its
        other equally-Dirac-free subalgebras, per points (2) and (3) above.
        Hence $F$ cannot be simply ``maximise $\dim \mathcal{D}$.''
  \item The spectral action functional
        $S[D] = \mathrm{Tr}\, f(D/\Lambda) + \langle \psi, D\psi \rangle$
        is the natural candidate, because the heat-kernel coefficients
        $a_n$ mix Dirac spectrum content with representation content of $A$.
        The question of whether the spectral action's value at the SM is
        actually a critical point in $A$ (not just in $D$) does not appear to
        have been settled in the literature as of May~2026.
  \item Khalkhali--Pagliaroli~\cite{khalkhali-pagliaroli} propose a moment-method
        / SDP bootstrap on Dirac ensembles, which provides numerical evidence
        for clustering of Dirac spectra near specific algebras. The bootstrap
        as currently formulated does not select $A_F$ uniquely either, but
        gives the closest existing tractable framework in which one could test
        a candidate $F$.
\end{enumerate}

We list these because they constitute a kind of negative specification: any
proposal for a Step~4 functional that fails any of (1)--(4) cannot be right,
and any proposal that passes all of them is worth taking seriously. We have
not produced such a proposal. This is, again, where the genuine open
question lives.

\bigskip\bigskip

\begin{thebibliography}{99}

\bibitem{wigner1960}
E.~P. Wigner, ``The Unreasonable Effectiveness of Mathematics in the Natural
Sciences,'' \emph{Comm.\ Pure Appl.\ Math.}\ \textbf{13} (1960) 1.

\bibitem{krajewski1998}
T.~Krajewski, ``Classification of Finite Spectral Triples,''
\emph{J.\ Geom.\ Phys.}\ \textbf{28} (1998) 1; arXiv:hep-th/9701081.

\bibitem{cc2007}
A.~H. Chamseddine and A.~Connes, ``Why the Standard Model,''
\emph{J.\ Geom.\ Phys.}\ \textbf{58} (2008) 38; arXiv:0706.3688.
Companion: \emph{Phys.\ Rev.\ Lett.}\ \textbf{99} (2007) 191601;
arXiv:0706.3690.

\bibitem{ccvs2013a}
A.~H. Chamseddine, A.~Connes, and W.~D. van Suijlekom, ``Inner Fluctuations
without the First Order Condition,'' \emph{J.\ Geom.\ Phys.}\ \textbf{73}
(2013) 222; arXiv:1304.7583.

\bibitem{ccvs2013b}
A.~H. Chamseddine, A.~Connes, and W.~D. van Suijlekom, ``Beyond the Spectral
Standard Model: Emergence of Pati--Salam Unification,'' \emph{JHEP} \textbf{11}
(2013) 132; arXiv:1304.8050.

\bibitem{boyle2018}
L.~Boyle and S.~Farnsworth, ``A new algebraic structure in the standard model
of particle physics,'' \emph{JHEP} \textbf{06} (2018) 071; arXiv:1604.00847.

\bibitem{ccvs2020}
A.~H. Chamseddine, A.~Connes, and W.~D. van Suijlekom, ``Entropy and the
Spectral Action,'' \emph{Comm.\ Math.\ Phys.}\ \textbf{373} (2020) 457;
arXiv:1809.02944.

\bibitem{connes-marcolli}
A.~Connes and M.~Marcolli, \emph{Noncommutative Geometry, Quantum Fields and
Motives}, AMS Colloquium Publications, vol.~55 (2008).

\bibitem{vansuijlekom-book}
W.~D. van Suijlekom, \emph{Noncommutative Geometry and Particle Physics},
Mathematical Physics Studies (Springer, 2015).

\bibitem{jureit-stephan}
J.-H.~Jureit and C.~A. Stephan, ``On a Classification of Irreducible
Almost-Commutative Geometries V,'' arXiv:0901.3214.

\bibitem{khalkhali-pagliaroli}
M.~Khalkhali and N.~Pagliaroli, ``Bootstrapping Noncommutative Geometry with
Dirac Ensembles,'' arXiv:2512.08694.

\bibitem{pdg2024}
S.~Navas \emph{et al.\ }(Particle Data Group),
``Review of Particle Physics,'' \emph{Phys.\ Rev.\ D} \textbf{110} (2024) 030001.

\bibitem{codata2022}
E.~Tiesinga, P.~J. Mohr, D.~B. Newell, and B.~N. Taylor, ``CODATA recommended
values of the fundamental physical constants: 2022,''
\emph{Rev.\ Mod.\ Phys.}\ \textbf{97} (2025) 025002.

\bibitem{strange-idea}
V.~Savytskyy with Claude (Anthropic), ``A Strange Idea, Followed Honestly:
The Generator Hypothesis After Every Stress Test We Could Apply,''
companion document, May 2026.

\bibitem{deflated-rewrite}
V.~Savytskyy and Claude (Anthropic), ``The Algebraic Core of Reality,
Honestly: A Maximally Deflated Rewrite,'' May 2026.

\end{thebibliography}

\end{document}
